% ===== tr-04 块2：第 4 章 pp.14--26；本片 pp.23--26（原文页 23--26）=====
% [原文 p.23]
带有第 2 节中定义的 McLean 度量与仿射结构。人们随后可以紧化 $B_{0,i}\subseteq B_i$（大概是通过取 $B_{0,i}$ 在特殊拉格朗日流（current）空间中的闭包；细节在此并不重要）。这给出一列度量空间 $(B_i,d_i)$，其中度量 $d_i$ 由 McLean 度量诱导。若把 McLean 度量标准化，使 $B_i$ 的直径与 $i$ 无关地保持常数，则我们可以期望 $(B_i,d_i)$ 收敛到一个紧度量空间 $(B_\infty,d_\infty)$。于是 SYZ 的极限形式如下：

\begin{conjecture}{4.5.5}
若 $(\mathcal{X}_{t_i},g_i)$ 收敛到 $(X_\infty,g_\infty)$，且 $(B_i,d_i)$ 对大的 $i$ 非空并收敛到 $(B_\infty,d_\infty)$，则 $B_\infty$ 与 $X_\infty$ 在相差一个缩放的意义下等距。此外，存在子空间 $B_{\infty,0}\subseteq B_\infty$，使得 $\Gamma:=B_\infty\setminus B_{\infty,0}$ 是 $B_\infty$ 中 Hausdorff 余维二的集合，且 $B_{\infty,0}$ 是一个 Monge--Ampère 流形，其 Monge--Ampère 度量在 $B_{\infty,0}$ 上诱导出 $d_\infty$。
\end{conjecture}
% [Tutor Note: 两个极限用的字母不同：流形一侧是 (X_∞, g_∞)（Riemann 度量），底一侧是 (B_∞, d_∞)（只给度量空间的结构）；结论是 B_∞ 与 X_∞ 等距（up to scaling）。另注意 Monge--Ampère 流形是 B_{∞,0}（去掉余维二集之后），不是整个 B_∞。]

这实质上是在说：当我们逼近极大幂零退化时，我们预期在 $\mathcal{X}_{t_i}$ 越来越大的子集上有一个特殊拉格朗日纤维化。此外，在极限中，Joyce 所提示的余维一判别轨迹收敛为余维二的判别轨迹，而 $B_{0,i}$ 上（不一定为 Monge--Ampère 的，见 [60]）Hessian 度量收敛为 Monge--Ampère 度量。

我在这里想说明的要点是：SYZ 猜想很可能只是「近似地」正确，人们必须看极限才有希望证明任何东西。另一方面，上面的猜想似乎很可能可以用目前已经掌握的技术达到。我仍然希望这个猜想会被证明，尽管还需要大量额外的工作。

我们如何用这一修改版的 SYZ 猜想来做镜像对称？本质上，我们将遵循以下步骤：
\begin{enumerate}
\item 我们从 Calabi--Yau 流形的一个极大幂零退化 $\mathcal{X}\to D$ 以及一个极化（polarization）的选取出发。这为每个 $t\in D\setminus0$ 给出一个 Kähler 类 $[\omega_t]\in H^2(\mathcal{X}_t,\R)$，它由 Ricci 平坦度量 $g_t$ 的 Kähler 形式 $\omega_t$ 代表。
\item 找出序列 $(\mathcal{X}_{t_i},r_ig_{t_i})$ 的 Gromov--Hausdorff 极限，其中 $t_i\to0$，而 $r_i$ 是使 $\mathcal{X}_{t_i}$ 的直径保持常数的缩放因子。若上述猜想成立，极限将是一个带奇点的仿射流形 $B$ 连同其上的 Monge--Ampère 度量。
\item 作一次 Legendre 变换，得到一个新的带奇点仿射流形 $\check B$，但度量相同。
\item 尝试构造 $X_\epsilon(\check B_0)$ 对小的 $\epsilon>0$ 的紧化，以得到一个复流形 $X_\epsilon(\check B)$。这将是镜像流形。
\end{enumerate}
% [Tutor Note: 步骤 (2) 中先缩放后取极限，度量是 r_i g_{t_i}（不是 g_{t_i}），r_i 只负责把直径固定；步骤 (4) 紧化的对象是 X_ε(ˇB_0)（而非 ˇX(B_0)），两者相差由 ε 控制的半平坦模型。]

正如我们将看到的，我们并不预期完成镜像对称需要步骤 (2) 与 (3) 的全部力量；某种识别底 $B$ 的办法就足够了。
% [原文 p.24]
不过，从理解 Ricci 平坦 Kähler 流形的微分几何的角度看，(2) 很有意思。

另一方面，步骤 (4) 至关重要，我们需要对最后这一步再多说几句。问题在于：虽然我们预期一般地应当可以构造辛流形 $\check X(B_0)$ 的辛紧化（并因此把镜像作为一个辛流形得到，三维情形见 [8]），但我们并不预期能把 $X_\epsilon(\check B_0)$ 作为复流形紧化。相反，预期的做法是先对 $X_\epsilon(\check B_0)$ 作一个小形变，然后才能紧化。此外，这个小形变在镜像对称中至关重要：正是这个小形变提供了 B-模型的瞬子修正。

由于最后这一项如此重要，让我们给它起个名字：

\begin{question}{4.5.6}[重构问题，版本 I]
给定一个带奇点的热带仿射流形 $B$，构造一个复流形 $X_\epsilon(B)$，它是 $X_\epsilon(B_0)$ 的某个小形变的紧化。
\end{question}
% [Tutor Note: 重构问题的要点是「先小形变、再紧化」这个顺序；X_ε(B_0) 本身一般无法作为复流形紧化，小形变不可省------它对应 B-模型的瞬子修正。]

我们稍后会在本文中回到这个问题。不过，我不希望在此进一步停留于 SYZ 猜想的微分几何版本。相反，我将接着描述上述讨论如何启发了由我与 Siebert 发展的、用于理解镜像对称的代数几何纲领，然后描述这一纲领带来的近期工作与想法。

\section{Gromov--Hausdorff 极限、代数退化与镜像对称}
我们现在有两种极限的概念：一方面是熟悉的代数几何概念，即圆盘上的退化族 $\mathcal{X}\to D$；另一方面是 Gromov--Hausdorff 极限。2000 年 Kontsevich 与 Soibelman 对这两者之间的联系有了一项重要的洞见（见 [53]）。本节我将粗略说明这是如何以及为何起作用的。

非常粗略地说，当 $t_i\to0$ 时 Gromov--Hausdorff 极限 $(\mathcal{X}_{t_i},g_{t_i})$------等价地，即设想的 SYZ 纤维化的底------应当在拓扑上与该奇异纤维 $\mathcal{X}_0$ 的对偶相交复形一致。更精确地说，在一个相对简单的情形，假设 $f:\mathcal{X}\to D$ 是（Mori 意义下）相对极小的且为正规交叉的，而 $\mathcal{X}_0$ 的不可约分支为 $X_1,\dots,X_m$。$\mathcal{X}_0$ 的对偶相交复形是以 $v_1,\dots,v_m$ 为顶点的单纯复形，当 $X_{i_0}\cap\cdots\cap X_{i_p}\ne\varnothing$ 时它含有一个单纯形 $\langle v_{i_0},\dots,v_{i_p}\rangle$。对偶相交复形应当在描述 SYZ 纤维化的底时起作用，这个想法也许最早由 Leung 与 Vafa 在 [56] 中提出。
% [Tutor Note: 对偶相交复形：顶点 v_1,…,v_m 对应不可约分支，单纯形 ⟨v_{i_0},…,v_{i_p}⟩ 存在 ⟺ 交集 X_{i_0}∩⋯∩X_{i_p} 非空（整体交集非空，不是两两相交）；文字层把 ⟨v_{i_0},…,v_{i_p}⟩ 抽成 ⟨vi0,...,vip⟩，多重下标 i_0…i_p 极易丢失，须按页图。]

让我们粗略解释为什么会这样，先看一族标准的退化椭圆曲线，其周期为 $1$ 与 $\frac{n}{2\pi i}\log t$，其中 $n$ 是正整数。这个
% [原文 p.25]
挖孔圆盘上的族可以通过在原点处添入一个 Kodaira 类型 $I_n$（$n$ 条有理曲线构成的循环）纤维而延拓为圆盘上的族。

取实数正序列 $t_i\to0$ 给出了一列形如 $X_{\epsilon_i}(B)$ 的椭圆曲线，其中 $B=\R/n\Z$ 且 $\epsilon_i=-\frac{2\pi}{\ln t_i}$。此外，$X_{\epsilon_i}(B)$ 上的度量在适当缩放之后来自 $B$ 上的常值 Hessian 度量。所以我们想解释 $B$ 如何与奇异纤维附近的几何相联系。为此，设 $X_1,\dots,X_n$ 是 $\mathcal{X}_0$ 的不可约分支；它们都是 $\PP^1$。设 $P_1,\dots,P_n$ 是 $\mathcal{X}_0$ 的奇点。

我们将考虑 $\mathcal{X}$ 中两类开集。对第一类，在 $X_i$ 上取坐标 $z$，使 $P_i$ 由 $z=0$ 给出、$P_{i+1}$ 由 $z=\infty$ 给出。设 $U_i\subseteq X_i$ 为开集 $\{z\mid\delta\le|z|\le1/\delta\}$，其中 $\delta$ 是某个固定的小量。则可以找到 $U_i$ 在 $\mathcal{X}$ 中的一个邻域 $\widetilde{U}_i$，使得 $\widetilde{U}_i$ 双全纯于 $U_i\times D_\rho$，其中 $\rho>0$ 充分小、$D_\rho$ 是 $\C$ 中半径为 $\rho$ 的圆盘，且 $f|_{\widetilde{U}_i}$ 是到 $D_\rho$ 的投影。
% [Tutor Note: U_i 的坐标条件 z=0 与 z=∞ 分别对应 P_i 与 P_{i+1}（下标是循环的，i+1 取模 n）；环带 {δ≤|z|≤1/δ} 的两端就是两个奇点附近，故在 𝒳 中它是个乘积区域。]

另一方面，每个 $P_i$ 在 $\mathcal{X}$ 中有一个邻域 $\widetilde{V}_i$ 双全纯于多圆盘
\[
\{(z_1,z_2)\in\C^2\mid |z_1|\le\delta',\ |z_2|\le\delta'\},
\]
在其上 $f$ 取为 $z_1z_2$ 的形式。

若 $\delta$ 与 $\delta'$ 选得恰当，则对充分接近零的 $t$，
\[
\{\widetilde{V}_i\cap\mathcal{X}_t\mid1\le i\le n\}\cup\{\widetilde{U}_i\cap\mathcal{X}_t\mid1\le i\le n\}
\]
构成 $\mathcal{X}_t$ 的一个开覆盖。现在这个开覆盖中的每个集合都可以写成 $X_\epsilon(U)$ 的形式，其中 $U$ 是一个一维（非紧）仿射流形，而 $\epsilon=-2\pi/\ln|t|$。若 $U$ 是开区间 $(a,b)\subseteq\R$，则 $X_\epsilon(U)$ 双全纯于环带
\[
\{z\in\C\mid e^{-2\pi b/\epsilon}\le|z|\le e^{-2\pi a/\epsilon}\},
\]
因为 $q=e^{2\pi i(x+iy)/\epsilon}$ 是 $X_\epsilon((a,b))$ 上的全纯坐标。于是
\[
\widetilde{U}_i\cap\mathcal{X}_t\cong X_\epsilon\Bigl(\Bigl(\frac{\epsilon\ln\delta}{2\pi},-\frac{\epsilon\ln\delta}{2\pi}\Bigr)\Bigr),
\]
其中 $\epsilon=-2\pi/\ln|t|$。当 $t\to0$ 时，区间 $(\epsilon\ln\delta/2\pi,-\epsilon\ln\delta/2\pi)$ 收缩为一点。所以当 $t\to0$ 时，按这种看待方式，$\widetilde{U}_i\cap\mathcal{X}_t$ 是 $\mathcal{X}_t$ 越来越小的开子集。这一论证提示每个不可约分支都应当与 $B$ 上的一个点相联系。
% [Tutor Note: 区间 (ε ln δ/2π, −ε ln δ/2π) 的左右端点符号相反（δ<1 时 ln δ<0），长度为 −ε ln δ/π；两个 ln δ 前的系数分别是 +ε/2π 与 −ε/2π，负号易丢，已在 600 dpi 下核对。]

现在看 $\widetilde{V}_i\cap\mathcal{X}_t$。它是
\[
\{(z_1,z_2)\in\C^2\mid |z_1|,|z_2|<\delta',\ z_1z_2=t\}\cong\{z\in\C\mid |t|/\delta'\le|z|\le\delta'\}
\]
\[
\cong X_\epsilon\Bigl(-\frac{\epsilon}{2\pi}\ln\delta',\ \frac{\epsilon}{2\pi}(\ln\delta'-\ln|t|)\Bigr),
\]
其中 $\epsilon=-2\pi/\ln|t|$。当 $t\to0$ 时，这个区间趋近于单位区间 $(0,1)$。于是开集 $\widetilde{V}_i\cap\mathcal{X}_t$ 最终成为 $\mathcal{X}_t$ 的一大部分。我们得到：对小 $t$，$\mathcal{X}_t$ 是形如 $X_\epsilon((i+\epsilon',i+1-\epsilon'))$ 的开集（即 $\widetilde{V}_i\cap\mathcal{X}_\epsilon$）与形如 $X_\epsilon((i-\epsilon'',i+\epsilon''))$ 的开集（即 $\widetilde{U}_i\cap\mathcal{X}_t$）之并，其中 $\epsilon',\epsilon''$ 充分小。这些开集应当（至少近似地）粘合起来给出
% [Tutor Note: 原文此句前半作 "(i.e., Ṽ_i ∩ 𝒳_ε)"，下标为 ε；而同句后半作 "Ṽ_i ∩ 𝒳_t"，按上下文（讨论 𝒳_t 的开覆盖）前半疑应为 𝒳_t。已在 600 dpi 下核实原文字母确为 ε，正文照录，存疑。]
% [原文 p.26]
$X_\epsilon(B)$。于是我们看到，$\mathcal{X}_0$ 的不可约分支似乎与 $B$ 上的点一致，而分支的交则与线段一致。这样我们便看到对偶相交复形浮现出来。

在下节开始严格结果之前，让我们再作一个观察。假设更一般地我们有一个 Calabi--Yau 流形的 \emph{Gorenstein 环面交叉}退化 $f:\mathcal{X}\to D$（见 [73]）。这意味着 $\mathcal{X}$ 的每一点 $x$ 都有一个邻域同构于某个仿射 Gorenstein（即典范类为 Cartier 除子）环面簇中的开集，而 $f$ 局部地由某个单项式给出，该单项式在每个余维一环面层上恰好以 1 阶消失。这是正规交叉概念的一个推广。非常粗略地说，上述论证提示中心纤维的每个不可约分支将对应于 Gromov--Hausdorff 极限的一个点。下面的习题说明，对 $\mathcal{X}_0$ 中一个零维层处的点 $x\in\mathcal{X}_0$，应当预期它对 $B$ 有什么样的贡献。
% [Tutor Note: Gorenstein 环面交叉退化的定义是「局部同构于仿射 Gorenstein 环面簇中的开集，且 f 由在每个余维一环面层上恰以 1 阶消失的单项式给出」；「toroidal crossings」不同于普通的 normal crossings（后者只是局部 z_1⋯z_k=0）。]

\begin{exercise}{4.6.1}
假设存在一点 $x\in\mathcal{X}_0$，它的一个邻域同构于某个仿射 Gorenstein 环面簇 $Y_x$ 的一个零维环面轨道的邻域。这样的仿射簇可如下指定。置 $M=\Z^n$、$M_{\R}=M\otimes_{\Z}\R$、$N=\operatorname{Hom}_{\Z}(M,\Z)$、$N_{\R}=N\otimes_{\Z}\R$，其中 $n=\dim\mathcal{X}_t$。则存在一个格多面体 $\sigma\subseteq M_{\R}$，
\[
C(\sigma):=\{(rm,r)\mid m\in\sigma,\ r\ge0\}\subseteq M_{\R}\oplus\R,
\]
$P:=C(\sigma)^\vee\cap(N\oplus\Z)$ 为锥 $C(\sigma)$ 的对偶所决定的幺半群，$Y_x=\operatorname{Spec}\C[P]$，最后 $f$ 等同于单项式 $z^{(0,1)}$。

现在取 $x$ 的一个如下形状的小邻域
\[
\widetilde{U}_\delta=\{y\in\operatorname{Spec}\C[P]\mid |z^p|<\delta\ \text{对所有}\ p\in P\}.
\]
这是一个开集，因为条件 $|z^p|<\delta$ 可以在 $P$ 的一个有限生成集上检验，只要 $\delta<1$。然后证明：对给定的 $t$、$|t|<1$ 且 $\epsilon=-2\pi/\log|t|$，若
\[
\sigma_t:=\Bigl\{m\in M_{\R}\Bigm|\langle p,(m,1)\rangle>\frac{\log\delta}{\log|t|}\ \text{对所有}\ p\in P\Bigr\},
\]
则
\[
f^{-1}(t)\cap\widetilde{U}_\delta\cong X_\epsilon(\sigma_t).
\]
注意
\[
\sigma:=\{m\in M_{\R}\mid\langle p,(m,1)\rangle\ge0\ \text{对所有}\ p\in P\},
\]
于是 $\sigma_t$ 是 $\sigma$ 的开子集，且当 $t\to0$ 时 $\sigma_t$ 收敛到 $\sigma$ 的内部。\proofend
\end{exercise}
% [Tutor Note: 习题 6.1 的配对是 ⟨p,(m,1)⟩（p ∈ P ⊂ N⊕Z，而 (m,1) ∈ M_R⊕R）；因 |t|<1 故 log|t|<0，而 δ<1 故 log δ<0，阈值 log δ/log|t| 为正。不等式取严格 > 才使 σ_t 成为开集；σ 里则取 ≥0。符号与不等号方向极易抄反。]

这个观察有希望为下一节的基本构造提供动机。
